\documentclass[11pt,a4paper]{article}

\usepackage[utf8]{inputenc}
\usepackage[T1]{fontenc}
\usepackage{lmodern}
\usepackage[margin=2.5cm]{geometry}
\usepackage{amsmath,amssymb}
\usepackage{booktabs}
\usepackage{tabularx}
\usepackage{enumitem}
\usepackage{xcolor}
\usepackage{fancyvrb}
\usepackage[most]{tcolorbox}
\usepackage[colorlinks=true,linkcolor=blue!60!black,urlcolor=blue!60!black]{hyperref}

\setlist{itemsep=2pt,topsep=4pt}
\setlength{\parindent}{0pt}
\setlength{\parskip}{6pt}
\setlength{\emergencystretch}{3em}

\newcommand{\ui}[1]{\textsf{\textbf{#1}}}          % menu entry / button
\newcommand{\file}[1]{\texttt{#1}}
\newcommand{\col}[1]{\texttt{#1}}                  % data column name

\newtcolorbox{note}{colback=blue!4,colframe=blue!40!black,boxrule=0.5pt,arc=2pt,
  left=6pt,right=6pt,top=4pt,bottom=4pt,fonttitle=\bfseries,title=Note}
\newtcolorbox{shell}{colback=gray!6,colframe=gray!50,boxrule=0.4pt,arc=2pt,
  left=6pt,right=6pt,top=2pt,bottom=2pt}

\title{\textbf{EXAFS Data Reduction}\\[4pt]
  \large User manual}
\date{\today}
\author{\url{https://github.com/tkroll123/exafs_data_reduction}}

\begin{document}
\maketitle

\begin{abstract}
\noindent EXAFS Data Reduction is a small web application for the first steps of
reducing X-ray absorption scans recorded with the SSRL \emph{EXAFS Data Collector}.
It reads the raw ASCII scan files, selects and cleans the fluorescence channels,
chooses the detector that carries the signal, calibrates the energy axis with a
reference foil, merges repeated scans and writes text files that document every
setting used. It runs on your own computer and is operated from a web browser.
The processing follows the XAS\_Worker scripts by Leland B.\ Gee, with all
hard-coded choices (columns, foil channels, thresholds) turned into settings.
\end{abstract}

\tableofcontents
\newpage

%=====================================================================
\section{Installation}
%=====================================================================

\subsection{Requirements}
\begin{itemize}
  \item \textbf{Python 3.9 or newer.} On macOS run \texttt{python3 --version} in Terminal; if
        Python is missing, install it from \url{https://www.python.org/downloads/}\\ or
        with Homebrew (\texttt{brew install python}). On Windows install Python from
        \url{https://www.python.org/downloads/} and tick \emph{Add python.exe to PATH}.
  \item \textbf{git} (optional, to download the code; you can also download a ZIP from GitHub).
  \item \textbf{A web browser.} The program was developed with Google Chrome; other current browsers
        should work as well.
  \item \textbf{Internet access the first time a page is opened}: the plotting library
        (Plotly) is loaded from \url{https://cdnjs.cloudflare.com}.
\end{itemize}

The Python packages used are Flask (web server), NumPy and SciPy. They are
installed into a \emph{virtual environment}, a folder \file{.venv} inside the
program folder, so nothing changes in your system Python.

\subsection{macOS and Linux}
\begin{shell}
\begin{Verbatim}
git clone https://github.com/tkroll123/exafs_data_reduction.git
cd exafs_data_reduction
python3 -m venv .venv
.venv/bin/pip install -r requirements.txt
\end{Verbatim}
\end{shell}
Without git: on the GitHub page click \ui{Code} $\to$ \ui{Download ZIP}, unpack it,
open a terminal in the unpacked folder and run the last two commands.

\subsection{Windows}
In PowerShell or the Command Prompt:
\begin{shell}
\begin{Verbatim}
git clone https://github.com/tkroll123/exafs_data_reduction.git
cd exafs_data_reduction
py -m venv .venv
.venv\Scripts\pip install -r requirements.txt
\end{Verbatim}
\end{shell}

\subsection{Updating}
\begin{shell}
\begin{Verbatim}
cd exafs_data_reduction
git pull
.venv/bin/pip install -r requirements.txt     # Windows: .venv\Scripts\pip ...
\end{Verbatim}
\end{shell}

%=====================================================================
\section{Starting and stopping}
%=====================================================================
From the program folder start the server:
\begin{shell}
\begin{Verbatim}
.venv/bin/python app.py          # Windows: .venv\Scripts\python app.py
\end{Verbatim}
\end{shell}
The terminal shows \texttt{Running on http://127.0.0.1:5050}. Open
\url{http://127.0.0.1:5050} in the browser. Keep the terminal open while you
work; stop the server with \texttt{Ctrl+C}.

\begin{note}
\begin{itemize}[leftmargin=*]
  \item The server only accepts connections from your own computer (\texttt{127.0.0.1}).
  \item Opened files are uploaded to the local server and kept in memory only. They are
        forgotten when the server stops or when you reload the page; your files on disk are never
        changed. After a restart, open the files again.
  \item If port 5050 is already in use, change \texttt{port=5050} in the last line of
        \file{app.py}.
\end{itemize}
\end{note}

%=====================================================================
\section{Overview of the window}
%=====================================================================
The window has a dark \textbf{menu bar on the left} with all settings and a
\textbf{plot area on the right} with six tabs.

\begin{center}
\begin{tabularx}{\linewidth}{@{}lX@{}}
\toprule
Menu section & Purpose \\
\midrule
\ui{Open files\ldots}        & Load one or more scan files. \\
\ui{Files}                   & List of open files: click a name to show it, tick the box to select it for
                               merging / downloading, $\times$ closes it. \ui{$\boxtimes$ all} ticks / unticks all,
                               \ui{clear} closes all. \\
\ui{Channels}                & Energy, I0 and I1 columns; auto-drop threshold. \\
\ui{Fluorescence channels}   & Which columns are fluorescence channels, which detector is used, and which channels
                               are dropped (clickable grid). \\
\ui{Energy calibration}      & Reference foil element, foil columns, method and its parameters. \\
\ui{Merge checked files}     & Combine the ticked scans. \\
\ui{Download processed data} & Save processed files. \\
\bottomrule
\end{tabularx}
\end{center}

\begin{center}
\begin{tabularx}{\linewidth}{@{}lX@{}}
\toprule
Tab & Shows \\
\midrule
\ui{File info}      & Header fields, which column was assigned to which role, raw header text. \\
\ui{Raw columns}    & Any columns against energy or against elapsed time (\ui{x = elapsed time}, e.g.\ to look
                      at I0 during the scan); optionally divided by I0, log axis, calibrated energy. \\
\ui{Channels / I0}  & Detector table with signal-to-noise; channel quality table (Section~\ref{sec:quality});
                      mean of every channel / I0 as bars (click to drop);
                      every used channel / I0 against energy. \\
\ui{Summed ratio}   & $\Sigma$ fluorescence / I0 on the calibrated energy axis, transmission $\ln(I_0/I_1)$,
                      statistics. \\
\ui{Calibration}    & Foil derivative near the edge with the point used for calibration, and the foil
                      $\mu(E)$. \\
\ui{Merged}         & Result of the last merge, the individual scans, their differences from the merge and a
                      table of the scans. \\
\bottomrule
\end{tabularx}
\end{center}

All plots are interactive (Plotly): drag to zoom, double-click to reset, click a
legend entry to hide a trace, the camera icon saves a PNG.

%=====================================================================
\section{Opening files}
%=====================================================================
Click \ui{Open files\ldots} and select one or several files (Shift/Cmd-click in the
dialog). The list is sorted by name with numbers in natural order
(\file{scan\_2} before \file{scan\_10}). Three formats are recognised automatically:

\begin{description}
  \item[SSRL EXAFS Data Collector] (e.g.\ \file{sample\_A.001}). The header is read up to the
        line \texttt{Data:}; the following lines up to the first empty line are the column names, then
        the numeric data. The header fields \texttt{PTS}/\texttt{COLS}, detector file, region file,
        scan parameters and comment lines are shown under \ui{File info}. The beamline is taken from the
        header line that starts with two digits and a word (e.g.\ \texttt{73 unfocused \ldots} $\to$ BL7-3)
        and shown as a badge in the file list.
  \item[Files written by this program] (\file{*\_processed.dat}, \file{*\_merged\_*.dat}). They are
        shown as stored, marked with a purple badge; see Section~\ref{sec:reopen}.
  \item[Generic ASCII] Any whitespace- or comma-separated numeric table. If the line above the
        numbers has as many words as there are columns, it is used as column names, otherwise the columns
        are called \col{col1}, \col{col2}, \ldots
\end{description}
Lines whose number of values does not match the number of columns are skipped and
counted (\ui{File info} $\to$ \emph{Rows read}). If the header gives a different
number of points than was read, a yellow warning is shown.

When the first file is opened, the foil element is guessed: the lowest tabulated
foil edge that lies at least 20\,eV above the start of the scan and below its end.

%=====================================================================
\section{Settings}
%=====================================================================

\subsection{How settings apply to several files}
\begin{itemize}
  \item \textbf{Column choices} (energy, I0, I1, foil numerator / denominator) start from each file's own
        guess. Once you choose a column in the menu, that choice is used for every open file that has a
        column of that name; files without it keep their own guess. \ui{default} (calibration section)
        goes back to the guessed foil columns.
  \item \textbf{Manual channel drops} are stored per file. \ui{Copy drops to all files} copies the drops
        of the file shown to all open files; \ui{Clear drops} clears them for the file shown.
  \item \textbf{All other settings} (threshold, detector, calibration element and method, \ldots) are
        global and apply to every file.
\end{itemize}
Every change is applied immediately and the plots of the current file are updated.

\subsection{Channels}
\begin{description}
  \item[Energy] Default \col{Achieved Energy}, else \col{Requested Energy}, \col{energy}, \col{eV}
        or the first column.
  \item[I0] Incident-intensity monitor, default \col{I0}.
  \item[I1] Transmission detector, default \col{I1}; ``none'' turns off the transmission curve.
  \item[Zero drop] A fluorescence channel with at least this many points equal to exactly zero is
        dropped automatically (default 5, as in XAS\_Worker; 0 switches this test off). See also the
        quality check in Section~\ref{sec:quality}.
\end{description}

\subsection{Fluorescence channels and detector}
\begin{description}
  \item[Pattern] \emph{auto} takes all columns named \col{SCA\textit{m}\_\textit{n}},
        \col{SCA\_\textit{n}} or \col{FF\textit{n}}. \emph{custom regex} lets you give a regular
        expression, e.g.\ \verb|^SCA1_\d+$|.
  \item[Detector] Channels are grouped by their name without the trailing number
        (\col{SCA1\_12} $\to$ \col{SCA1}). Often only one detector carries the fluorescence signal. With
        \emph{auto} the detector with the highest edge-jump signal-to-noise ratio (below) is used and its
        name is shown in the drop-down; you can also choose one detector or \emph{all detectors}. With only
        one detector the menu is inactive.
  \item[Distortion $\sigma$] Limit of the distortion test (Section~\ref{sec:quality}), default 15;
        0 switches it off.
  \item[Outlier limit] Strictness of the noise and spike tests, default 3; 0 switches them off.
        The tick boxes \ui{distortion}, \ui{noise} and \ui{spikes} choose which tests are used
        (default: distortion and noise).
\end{description}

The \textbf{channel grid} shows each channel with a colour:
\textcolor{green!50!black}{\textbf{green}} = used,
\textcolor{red!70!black}{\textbf{red}} = dropped automatically,
\textcolor{gray}{\textbf{grey, crossed out}} = dropped by you,
dark = belongs to a detector that is not used,
green with a \textcolor{orange}{\textbf{yellow frame}} = kept by you although the automatic check would drop it.
Hovering shows why a channel was dropped and its quality values. Clicking a channel
(or its bar in the \ui{Channels / I0} tab) works as follows:
used $\to$ manual drop $\to$ used; automatically dropped $\to$ kept anyway $\to$ automatically dropped.
Manual drops and keeps are stored per file (\ui{Copy drops to all files} copies both).

\subsection{Automatic channel quality check}\label{sec:quality}
Every channel is compared with a \emph{reference} spectrum made from the clean channels of the
same detector, which all see the same sample. With $r_c = F_c/I_0$ and $E_0$ the tabulated foil edge:
\begin{enumerate}
  \item pre-edge level $p_c$ = mean of $r_c$ at $E_0-150\ldots E_0-20$\,eV, post-edge level $q_c$ at
        $E_0+50\ldots E_0+200$\,eV, edge jump $j_c = q_c-p_c$, normalized spectrum $n_c=(r_c-p_c)/j_c$;
  \item first reference: at each energy the mean of the \emph{tightest third} of the $n_c$ (``shorth'').
        The clean channels agree closely and form this tight group, even when more than half of
        the channels are disturbed at that energy (a plain median would be pulled along);
  \item each channel is fitted robustly as $r_c \approx b_c + a_c\cdot\text{reference}$ (background
        $b_c$, signal $a_c$); the residual $(r_c-\text{model})/a_c$ is what is special about the channel;
  \item the reference is rebuilt three times as the median of the 40\,\% of channels with the smallest
        distortion (below).
\end{enumerate}
From the residual these numbers are computed (shown in the \ui{Channel quality} table of the
\ui{Channels / I0} tab and in the tooltip of each channel):
\begin{center}
\begin{tabularx}{\linewidth}{@{}lX@{}}
\toprule
Value & Meaning and rule\\
\midrule
Edge jump    & $j_c$ in \% of the median jump. Below 30\,\% the channel is \textbf{dead} and always dropped.\\
Distortion   & largest deviation of the 5-point running mean of the residual, divided by the noise of
               that running mean (from the channel's own point-to-point scatter), i.e.\ in $\sigma$.
               The table also gives its size (\% of the edge jump) and energy.
               Dropped when above \ui{Distortion $\sigma$} (default 15). This catches sudden drops
               or steps that do not average out when scans are merged, in particular \textbf{dead-time
               losses from diffraction}: when a Bragg peak of the sample hits some detector elements,
               their total count rate rises, fewer counts fall into the fluorescence window, and
               $F_c/I_0$ drops suddenly and recovers slowly as the peak moves on with energy.
               The limit is absolute, because often most channels are hit at the same energy.\\
Noise        & robust (median-absolute-deviation) scatter of the residual divided by the
               counting-statistics expectation $\sqrt{\text{counts}}$; 1 = ideal. Dropped when above
               median $+\,k\cdot$MAD of the live channels, $k$ = \ui{Outlier limit} (default 3).\\
Spikes       & number of points deviating by more than $6\sigma$ (Poisson); same relative rule.
               Off by default.\\
\bottomrule
\end{tabularx}
\end{center}
The tick boxes \ui{distortion}, \ui{noise} and \ui{spikes} switch the tests on or off.

\paragraph{Why these defaults.} The rules were tested on 22 Fe K-edge scans (24-element detector)
for which the channels had been selected by hand.
\begin{itemize}
  \item Many scans show sudden drops of 5--30\,\% in a group of channels (e.g.\ \file{acid\_015}
        at 7600\,eV). The distortion test at 15$\sigma$ removes them: 144 channels in total, of which
        131 had also been removed by hand; in \file{acid\_015} it drops exactly the 13 disturbed channels
        plus the dead one.
  \item The hand selections removed 5--19 of 24 channels; the summed spectrum with the automatic
        selection is still about 10\,\% less noisy, because channels that are only slightly
        different were kept.
  \item The spike test removed many more channels, increased the noise of the sum and did not reduce
        the number of spikes in it (these are glitches common to all channels), so it is off by default.
\end{itemize}
A lower \ui{Distortion $\sigma$} (e.g.\ 10) also removes channels with smaller dips; clean channels
typically lie at 5--10$\sigma$. The quality check needs a foil element and a scan covering
$E_0-150\ldots E_0+200$\,eV; otherwise only the zero-count test is used and a note is shown.

\paragraph{Detector score.} For each detector the used channels are summed and
divided by I0, giving $y(E)$. With $E_0$ the tabulated foil energy:
\begin{align*}
  \text{jump}  &= \langle y\rangle_{E_0+30\ldots E_0+100\,\mathrm{eV}} - \langle y\rangle_{E_0-80\ldots E_0-15\,\mathrm{eV}},\\
  \text{noise} &= \operatorname{std}\bigl(y_{i+1}-y_i\bigr)_{\text{pre-edge}}\big/\sqrt{2},\qquad
  \text{S/N} = \text{jump}/\text{noise}.
\end{align*}
If the scan does not cover both ranges (at least 4 pre-edge and 3 post-edge points),
the detector with the highest mean ratio is used instead.

\subsection{Energy calibration}
The energy axis is shifted so that the edge of the reference foil, measured at the
same time as the sample, lands on its tabulated value:
\[
  E_\text{cal} = E_\text{raw} + \Delta, \qquad \Delta = E_\text{ref} - E_\text{found}.
\]
\begin{description}
  \item[Calibrate with reference foil] Off: energies are not shifted.
  \item[Element] Foil element, which sets $E_\text{ref}$ (Table~\ref{tab:foils}).
  \item[Foil num.\ / denom.] Columns of the intensities before and after the foil,
        $\mu_\text{foil} = \ln(\text{num}/\text{denom})$. Default \col{I1}/\col{I2}; for files from BL7-3
        \col{I0}/\col{Scatter}. The note below the menus tells which default was used, or warns (yellow) when the
        beamline default columns are missing.
  \item[Method] see below.
  \item[Smooth FWHM] Width (eV) of the Gaussian smoothing of $\mu$ (0 = off). Default 1.0\,eV.
  \item[Spline step] Grid spacing (eV) on which the spline derivative is evaluated (0 = no spline).
        Default 0.01\,eV.
  \item[Search $\pm$ eV] The maximum of the derivative is searched within $E_\text{ref}\pm$ this value.
        Default 10\,eV.
\end{description}

\paragraph{Method ``ln''} (default):
\begin{enumerate}
  \item $\mu(E)=\ln(\text{num}/\text{denom})$, sorted by energy; invalid points are removed and duplicate
        energies averaged.
  \item Gaussian smoothing: every point becomes the Gaussian-weighted mean of its neighbours within
        $\pm4\sigma$, $\sigma = \text{FWHM}/2.3548$ (works on non-uniform energy grids).
  \item A cubic spline through the smoothed $\mu$ (not-a-knot end conditions, identical to MATLAB
        \texttt{spline}).
  \item The spline derivative $d\mu/dE$ is evaluated every \emph{spline step} eV inside the search window;
        its maximum is $E_\text{found}$. Without spline, finite differences on the data points are used.
\end{enumerate}
If the maximum lies at the border of the search window, a warning is shown: the
element or the foil columns are probably wrong, or the window is too small.

\paragraph{Method ``legacy''} reproduces XAS\_Worker exactly: $\mu = \log_{10}(r - \min r)$
with $r=\text{num}/\text{denom}$, derivative $\mu_{i}-\mu_{i-2}$, maximum within
$(E_\text{ref}-w,\,E_\text{ref}+w]$, no smoothing or spline.

\paragraph{The Calibration tab} shows the derivative in the uncalibrated energy
scale: grey points = raw derivative, blue points = derivative of the smoothed
$\mu$, blue line = spline derivative. The shaded band is the search window, the
dotted grey line the tabulated energy (target), the red dashed line and circle the
maximum that was used. The lower plot shows the foil $\mu(E)$.

\begin{table}[h]
\centering
\caption{Tabulated foil K-edge energies (eV). More can be added to \texttt{FOILS} in \file{xas\_core.py}.}
\label{tab:foils}
\begin{tabular}{@{}cccccccccc@{}}
\toprule
Sc & Ti & V & Cr & Mn & Fe & Co & Ni & Cu & Zn\\
\midrule
4492.8 & 4966.4 & 5465.1 & 5989.2 & 6539.0 & 7111.2 & 7708.9 & 8332.8 & 8978.9 & 9658.6\\
\bottomrule
\end{tabular}
\end{table}

%=====================================================================
\section{Results of one scan}
%=====================================================================
\[
  \text{SumRatio}(E_\text{cal}) = \frac{\sum_{c\,\in\,\text{used}} F_c}{I_0}\;\langle I_0\rangle,
  \qquad
  \text{transmission} = \ln\frac{I_0}{I_1}.
\]
$\langle I_0\rangle$ is the mean I0 of the scan (written to the header as \texttt{I0\_mean}; values
$\le 0$ are ignored). Multiplying by it keeps SumRatio in units of counts at the average beam
intensity: a scan taken with a weaker beam gives a correspondingly smaller SumRatio, so the
weight of a scan reflects the incident flux it actually received (for a sample of constant
concentration).
The \ui{Summed ratio} tab shows both curves, the calibration shift, the used and
dropped channels, and the statistics written to the output file (mean, minimum and
maximum of SumRatio, I0, I1; energy range; shift).

%=====================================================================
\section{Merging scans}
%=====================================================================
Tick the scans to combine in the file list (at least two) and click
\ui{Merge checked}. Each scan is first processed with its own settings
(columns, detector, drops, calibration). The scans are then combined by
\textbf{adding counts}, as if all good channels of all scans were one large detector:
\[
  \text{merged}(E) = \frac{\sum_\text{scans}\sum_{c\,\in\,\text{used}} F_c(E)}{\sum_\text{scans} I_0(E)}\;\langle I_0\rangle,
  \qquad
  \text{transmission} = \ln\frac{\sum I_0}{\sum I_1}.
\]
No scan is rescaled; each contributes in proportion to its counts (shown as
``Share of counts'' in the table). $\langle I_0\rangle$ is the mean I0 of one point of one scan,
averaged over all merged scans and points (\texttt{I0\_mean} in the header), so the merged
curve is on the same scale as the single scans.

\begin{description}
  \item[Energy grid: interpolate to first scan] (default) Every scan is linearly interpolated onto the
        calibrated energies of the first ticked scan (in file-list order). There is no extrapolation, and
        only grid points covered by \emph{all} scans are kept: because counts are added, a point where one
        scan is missing would get a different mix of scans and jump (e.g.\ the last point, when the
        calibration shifts differ by a few hundredths of an eV). The left-out points are listed in the
        \ui{Merged} tab and in the header (\texttt{points\_left\_out}); usually it is the first and/or last point.
  \item[Energy grid: point-by-point] Points are combined by index, the energy is the mean of the
        calibrated energies (XAS\_Worker default). All scans must have the same number of points.
\end{description}
The \ui{Merged} tab keeps the result of the last merge; after changing settings,
click \ui{Merge checked} again. \ui{Download merged} recomputes the merge with the
current settings and saves it.

Already processed or merged files (Section~\ref{sec:reopen}) can be ticked and merged
together with raw scans.

%=====================================================================
\section{Downloading results}
%=====================================================================
\begin{description}
  \item[Download current file] saves the scan shown as \file{<name>\_processed.dat}.
  \item[Download selected (N)] saves every ticked file, one after the other, each as its own
        \file{\_processed.dat}. N is the number of ticked files.
  \item[Download merged] saves \file{<common prefix>\_merged\_<N>scans.dat}.
  \item[Overhead s/pt] time per point spent moving motors and reading out, which the scan file does
        not record (default 1.7\,s). It is used only for the column \col{Elapsed\_s}. The default was
        measured at BL7-3: consecutive scans of 506 points with 1195\,s counting time started
        2030--2070\,s apart, i.e.\ about 1.7\,s per point.
\end{description}
The name follows XAS\_Worker: \file{sample\_A.001} $\to$ \file{sample\_A001\_processed.dat}.
Files are saved to the browser's download folder. If Chrome's setting \emph{Ask where
to save each file before downloading} is on, a dialog appears for every file; the
first time several files are downloaded, Chrome may ask to allow multiple downloads.

\begin{note}
Processed or merged files that were opened again cannot be exported again as
\file{\_processed.dat} (they are already processed); they are reported with an
error message and the other files are still saved.
\end{note}

\subsection{Format of a processed file}
Plain text. Header lines start with \texttt{\#} and are grouped in sections:
\begin{description}[font=\ttfamily]
  \item[[source]] original file, format, collection date, beamline, detector / region files, scan
        parameters, comments, number of points.
  \item[[channels]] energy / I0 / I1 columns, detector mode and detector used, a line per detector with
        channels used, edge jump, noise and S/N, then the lists of used, automatically dropped (with the
        rules and the reason for each channel), kept-by-you, manually dropped and switched-off channels.
  \item[[energy calibration]] element, reference energy, foil columns, method, smoothing, spline step,
        search window, edge found and shift (or the error).
  \item[[statistics]] mean/min/max of SumRatio, I0, I1, the energy range and \texttt{I0\_mean}.
  \item[[time]] \texttt{scan\_start} (date and time from the scan header), the count-time columns used,
        the overhead per point and the estimated scan duration.
  \item[[data]] the column names, then the data, sorted by energy.
\end{description}
\begin{center}
\begin{tabularx}{\linewidth}{@{}lX@{}}
\toprule
Column & Meaning\\
\midrule
\col{Energy\_cal} & calibrated energy (eV) $= E_\text{raw} + \Delta$\\
\col{Energy\_raw} & energy as recorded (eV)\\
\col{SumRatio}    & $\sum$ used channels / I0 $\times$ \texttt{I0\_mean}\\
\col{sum\_of\_used\_SCA1} & sum of the raw counts of the used channels; named after the detector in use
                    (e.g.\ \col{sum\_of\_used\_SCA2}, or \col{sum\_of\_used\_SCA1\_SCA2} with \emph{all detectors})\\
\col{I0}          & I0 counts\\
\col{Scatter}     & raw counts of the foil denominator column (and of the numerator, if it is not I0)\\
\col{ln(I0/Scatter)} & reference foil absorbance $\ln(\text{num}/\text{denom})$ used for the calibration;
                    the name follows the foil columns chosen (e.g.\ \col{ln(I1/I2)})\\
\col{Count\_time\_s} & counting time of the point (s), from \col{Real time clock}\\
\col{Sum\_count\_time\_s} & cumulative counting time since the scan start (s), from \col{Sum\_RTC}\\
\col{Elapsed\_s}  & time from \texttt{scan\_start} to the middle of the point (s):
                    $\text{Sum\_count\_time} - \text{Count\_time}/2 + \text{overhead}\times\text{point number}$\\
\bottomrule
\end{tabularx}
\end{center}
Example header excerpt:
\begin{shell}
\begin{Verbatim}[fontsize=\small]
# EXAFS Data Reduction (Flask) - processed XAS scan
# [source]
# file: CoIIIoximeS1_s2_125_A.001
# beamline: BL7-3   (73 unfocused  0.000 0 X 1.920000 20000)
# [channels]
# detector_used: SCA1
# detector SCA1: used 30/30 channels, edge_jump=1.10234,
#                noise=0.00053992, S/N=2042  <- selected
# channels_used (30): SCA1_1 SCA1_2 ... SCA1_30
# auto_dropped (2): SCA1_7 SCA1_11
# auto_drop_reasons: SCA1_7=dead SCA1_11=distortion,noise
# [energy calibration]
# calibrated: yes
# element: Co
# reference_energy_eV: 7708.9
# foil: ln(I1/I2)
# edge_found_eV: 7707.6200
# energy_shift_eV: +1.2800   (calibrated = raw + shift)
# [data]
# [time]
# scan_start: Mon Nov 23 22:49:10 2026   (from the file header)
# overhead_per_point_s: 1.7
# [data]
# Energy_cal  Energy_raw  SumRatio  sum_of_used_SCA1  I0  Scatter  ln(I0/Scatter)
#   Count_time_s  Sum_count_time_s  Elapsed_s
\end{Verbatim}
\end{shell}

\subsection{Format of a merged file}
The header lists the merge mode, every scan with detector, number of channels, shift,
share of counts and dropped channels, per-scan details (comments, columns,
calibration), and a summary of the shifts. Data columns:
\col{Energy}, \col{SumRatio} (= \col{Sum\_counts}/\col{Sum\_I0}\,$\times$\,\texttt{I0\_mean}), \col{Sum\_counts},
\col{Sum\_I0}, \col{N\_scans} (scans contributing to that point) and, if all scans
have foil columns, the summed foil denominator (e.g.\ \col{Sum\_Scatter}) and the reference
of the summed foil signals, e.g.\ \col{ln(sumI0/sumScatter)}. Time columns are not merged.
Processed files written by earlier versions (with \col{ln(I0/I1)}) can still be opened and merged.

\subsection{Opening processed files again}\label{sec:reopen}
Files written by this program are recognised by their first line and shown as
stored, with their settings taken from the header (\ui{File info}, \ui{Summed ratio},
\ui{Calibration}). Menu settings do not apply to them, since the energies are already
calibrated. They can be merged with other files; their counts are recovered exactly
from the \col{sum\_of\_used\_\ldots} or \col{Sum\_counts} column (older files: SumRatio\,$\times$\,I0).

%=====================================================================
\section{Typical workflow}
%=====================================================================
\begin{enumerate}
  \item Start the server and open all scans of one sample.
  \item Look at \ui{File info} for the first scan: are energy, I0, I1, foil and fluorescence columns
        correct? Change them in the menu if not.
  \item Check the element and \ui{Calibration}: the red marker should sit on the main peak of the foil
        derivative close to the dotted target line.
  \item In \ui{Channels / I0} check which detector is used and look at the \ui{Channel quality} table.
        Normally the automatic check is enough; drop further channels, or keep an automatically dropped one,
        by clicking. Dropping many good channels makes the sum noisier, so only drop channels that are
        clearly worse than the others.
  \item Click through the other files and check calibration and channels for each.
  \item Tick the good scans, \ui{Merge checked}, inspect the \ui{Merged} tab (differences should be
        noise only), then \ui{Download merged} and/or \ui{Download selected}.
\end{enumerate}

%=====================================================================
\section{Troubleshooting}
%=====================================================================
\begin{description}
  \item[The page does not open.] Is the server still running in the terminal? Is the address
        \url{http://127.0.0.1:5050}?
  \item[Plots are empty, buttons do nothing.] The plotting library could not be loaded: check the internet
        connection and reload the page.
  \item[``Address already in use''.] Another program (or an older copy of this server) uses port 5050.
        Stop it, or change the port in \file{app.py}.
  \item[``unknown scan (re-open the file)''.] The server was restarted; open the files again.
  \item[``No foil data within $\pm$\ldots''.] The chosen element's edge is not inside the scan; choose the
        right element.
  \item[Calibration warning ``maximum at the edge of the search window''.] Wrong element or foil columns,
        a noisy foil, or the search window is too narrow.
  \item[No fluorescence channels found.] The column names do not follow \col{SCA\ldots\_n} / \col{FFn};
        use \emph{custom regex}.
  \item[Point-by-point merge fails.] Scans have different numbers of points; use
        \emph{interpolate to first scan}.
\end{description}

%=====================================================================
\section{For developers}
%=====================================================================
\begin{center}
\begin{tabularx}{\linewidth}{@{}lX@{}}
\toprule
File & Contents\\
\midrule
\file{app.py}               & Flask server, scans kept in a dictionary in memory.\\
\file{xas\_core.py}         & \texttt{parse\_ssrl}, \texttt{guess\_roles}, \texttt{calibrate}, \texttt{process},
                              \texttt{merge}, \texttt{export\_text}, \texttt{export\_merged\_text}.\\
\file{templates/index.html} & User interface (HTML, CSS, JavaScript, Plotly 2.34).\\
\bottomrule
\end{tabularx}
\end{center}

HTTP endpoints (all JSON unless noted):
\begin{center}
\begin{tabular}{@{}lll@{}}
\toprule
Method & Path & Purpose\\
\midrule
POST   & \texttt{/api/open}           & upload files (multipart field \texttt{files})\\
GET    & \texttt{/api/scan/<id>}      & header, columns and data of an opened scan\\
POST   & \texttt{/api/process/<id>}   & process with the settings in the body\\
POST   & \texttt{/api/export}         & processed file (text download)\\
POST   & \texttt{/api/merge}          & merge \texttt{\{items:[\{id,cfg\}], mode\}}\\
POST   & \texttt{/api/merge/export}   & merged file (text download)\\
DELETE & \texttt{/api/scan/<id>}      & close a scan\\
\bottomrule
\end{tabular}
\end{center}

The functions in \file{xas\_core.py} can also be used directly from Python, e.g.
\begin{shell}
\begin{Verbatim}[fontsize=\small]
import xas_core as xc
scan = xc.parse_ssrl(open('sample_A.001').read(), 'sample_A.001')
r = xc.guess_roles(scan['columns'], scan['meta'].get('beamline_guess'))
cfg = dict(energy=r['energy'], i0=r['i0'], i1=r['i1'],
           fluor_channels=r['fluor_channels'], manual_drops=[],
           auto_drop_threshold=5, calibrate=True, element='Co',
           foil_num=r['foil_num'], foil_den=r['foil_den'], foil_method='ln',
           detector='auto', smooth_fwhm=1.0, spline_step=0.01, cal_window=10)
res = xc.process(scan, cfg)
open('sample_A001_processed.dat', 'w').write(xc.export_text(scan, cfg, res))
\end{Verbatim}
\end{shell}

%=====================================================================
\section*{Acknowledgements}
%=====================================================================
The processing steps (automatic channel dropping, summed fluorescence ratio, foil
calibration, merging) follow XAS\_Worker by Leland B.\ Gee (MIT License, \copyright\ 2023).

\end{document}
